\chapter{Proofs a Computer Can Check}
\label{ch:lean}

\begin{center}\itshape
A referee reads a proof once, tired, on a train.\\
A proof assistant reads it every time, and never gets tired.\\
It also never understands it. That part is still yours.
\end{center}

In 1998 Thomas Hales announced a proof of the Kepler conjecture: no packing of equal balls in space is denser than the greengrocer's stack of oranges. The proof combined
three hundred pages of mathematics with billions of computer-checked inequalities. The \emph{Annals of Mathematics} sent it to a panel of referees, who worked on it for four
years and then reported that they were ``99\% certain'' it was correct, but could not certify the computations. The paper appeared in 2005. Hales's answer was not to argue; it was
to start a project, Flyspeck, to check the whole proof line by line in a proof assistant. It took a team of twenty-two people eleven years, and the result, ``A formal proof of the
Kepler conjecture'', appeared in 2017. Nobody is 99\% certain any more.

This chapter explains what such a check is, shows a few short proofs written for the proof assistant Lean~4 and its library Mathlib, and proves the theorem that makes the
whole enterprise possible: a proof is a program, and checking a proof is checking that a program has the right type.

\section{Four stories}

\begin{itemize}
\item \emph{The four colour theorem.} Appel and Haken's 1976 proof needed a computer to check over a thousand configurations, and for decades it was accepted with a shrug. In 2005
Georges Gonthier, with Benjamin Werner, gave a complete formal proof in the proof assistant Coq; his account in the \emph{Notices of the AMS} (2008) is a model of how to explain
such work to mathematicians.
\item \emph{Flyspeck}, above: the formal proof of the Kepler conjecture, completed in 2014 and published in 2017.
\item \emph{The Liquid Tensor Experiment.} In December 2020 Peter Scholze posted a challenge on Kevin Buzzard's blog: formalise a theorem of his with Dustin Clausen on liquid
vector spaces, of which he wrote ``I think this may be my most important theorem to date''. He was not sure the proof was right in every detail. A team led by Johan Commelin and
Adam Topaz finished the formalisation in Lean on 14 July 2022. On the way, the formalisers found simplifications of parts of the argument.
\item \emph{The polynomial Freiman--Ruzsa conjecture.} In November 2023 Gowers, Green, Manners and Tao posted a proof of this conjecture of additive combinatorics. Tao immediately
organised an open formalisation in Lean; about three weeks later, in early December 2023, it was complete.
\end{itemize}

Since October 2024 Buzzard leads a five-year project, funded by the UK research council EPSRC, to formalise a proof of Fermat's Last Theorem, with the stated aim of reducing it to results
known in the 1980s. The pattern is clear: formal proof has moved from checking old theorems to checking new ones as they are written.

\section{What a proof assistant does}

A proof assistant is a programming language whose types can express mathematical statements. A statement is a type; a proof of the statement is a program (a \emph{term}) of that type;
and the assistant's \emph{kernel}, a small program that one can read and trust, checks that the term has the type. Everything else (the tactics that help you build the term, the
automation, the library) can be as clever and as buggy as it likes: if it produces a term, the kernel checks it.

Here is the simplest possible example in Lean~4. The text after \texttt{--} is a comment.
\begin{verbatim}
-- a statement, and a proof of it
theorem two_add_two : 2 + 2 = 4 := rfl
\end{verbatim}
The statement \texttt{2 + 2 = 4} is a type; \texttt{rfl} (``reflexivity'') is a term of type \texttt{a = a}, and the kernel accepts it because both sides compute to the same natural number.
Implication is literally a function type: a proof of $A\to B$ is a function sending proofs of $A$ to proofs of $B$.
\begin{verbatim}
theorem mp (A B : Prop)
    (ha : A) (hab : A -> B) : B :=
  hab ha

theorem k (A B : Prop) : A -> B -> A :=
  fun a _ => a
\end{verbatim}
Modus ponens is function application, and the tautology $A\to(B\to A)$ is the function that takes $a$, ignores the second argument, and returns $a$. Section~\ref{sec:curryhoward} turns this
observation into a theorem.

\emph{Dependent types in one paragraph.} In Lean, a type may depend on a value: \texttt{Fin n} is the type of natural numbers below \texttt{n}, so a function \texttt{(n : Nat) -> Fin (n+1)}
returns a value whose type depends on its input. This is exactly what a universal statement needs. A proof of $\forall n,\ P(n)$ is a function that sends each $n$ to a proof of $P(n)$,
whose type depends on $n$. Existential statements are pairs: a witness and a proof that it works.

\emph{Tactics.} Writing terms by hand is like writing machine code. Instead one writes a proof script of \emph{tactics}, each of which transforms the current goal: \texttt{intro},
\texttt{apply}, \texttt{induction}, \texttt{simp} (rewriting with known equalities), \texttt{ring} (proving identities in commutative rings), \texttt{omega} (linear arithmetic over the integers),
\texttt{decide} (deciding a statement about finitely many cases by computation). Here is induction:
\begin{verbatim}
import Mathlib

def oddSum : Nat -> Nat
  | 0 => 0
  | n + 1 => oddSum n + (2 * n + 1)

theorem oddSum_eq (n : Nat) :
    oddSum n = n * n := by
  induction n with
  | zero => rfl
  | succ n ih =>
    rw [oddSum, ih]
    ring
\end{verbatim}
After \texttt{rw} (rewrite) the goal is $n\cdot n+(2n+1)=(n+1)(n+1)$, which \texttt{ring} proves. A human reader sees the familiar proof that $1+3+\dots+(2n-1)=n^2$; the kernel sees a term
built by the recursor of the natural numbers.

\section{Two theorems of the series, formalised}

Euclid's theorem, in the form ``above every $n$ there is a prime'', with the proof of Volume~1: a prime factor of $n!+1$ cannot be at most $n$. The function \texttt{Nat.minFac} returns the
smallest prime factor of a number different from $1$.
\begin{verbatim}
open Nat in
theorem euclid (n : Nat) :
    ∃ p, n ≤ p ∧ p.Prime := by
  let p := minFac (n ! + 1)
  have f1 : n ! + 1 ≠ 1 := by
    have := factorial_pos n
    omega
  have pp : p.Prime := minFac_prime f1
  refine ⟨p, ?_, pp⟩
  by_contra h
  push_neg at h
  have h1 : p ∣ n ! :=
    dvd_factorial (minFac_pos _) h.le
  have h2 : p ∣ n ! + 1 := minFac_dvd _
  have h3 : p ∣ 1 :=
    (Nat.dvd_add_right h1).mp h2
  exact pp.not_dvd_one h3
\end{verbatim}
Read it line by line against the proof you know: \texttt{f1} says $n!+1\ne1$ (because $n!>0$), so it has a smallest prime factor $p$; if $p<n$ (the hypothesis \texttt{h} after \texttt{push\_neg}),
then $p$ divides $n!$ and $n!+1$, hence $1$, which no prime does. Every step a human would skip (that $n!>0$, that $p\le n$ follows from $p<n$) has to be said, but only once:
the library then knows it for ever.

The obstruction of Proposition~\ref{prop:mod9} is a computation in $\Z/9\Z$, and computations are what \texttt{decide} is for:
\begin{verbatim}
theorem cubes_mod_nine (x : ZMod 9) :
    x^3 = 0 ∨ x^3 = 1 ∨ x^3 = -1 := by
  revert x
  decide
\end{verbatim}
The kernel checks the nine cases by evaluation. Turning this into the statement about integers is Exercise~\ref{ex:leanmod9}.

\begin{remark}[about the listings]\label{rem:leanlistings}
The listings of this chapter are written for Lean~4 with a recent Mathlib, and use only standard tactics and long-established library names. They were not compiled in the
environment in which this edition was produced, and library names do change over time: the first exercise of the chapter is to paste them into the online Lean editor
(\url{https://live.lean-lang.org}), which loads Mathlib, and repair whatever the current version objects to. Doing so is also the quickest way to learn how a proof assistant talks back.
\end{remark}

\section{Proofs are programs}
\label{sec:curryhoward}

Why should a type checker be able to check proofs at all? The answer, found independently by Haskell Curry (1934, for combinators) and William Howard (in notes written in 1969 and
published in 1980), is that proofs in a natural deduction system and programs in the typed $\lambda$-calculus are the same objects. We prove it for the simplest fragment of logic, where
the only connective is implication; the full correspondence (with $\wedge$, $\vee$, $\forall$, $\exists$) is the same idea with more cases.

\begin{definition}\label{def:nd}
\emph{Formulas} are built from atoms $p,q,\dots$ by $A\to B$. A \emph{context} is a finite list $\Gamma=(A_1,\dots,A_k)$ of formulas. \emph{Derivations} of judgements $\Gamma\vdash A$ are built by three rules:
\begin{enumerate}[label=(\roman*)]
\item (hypothesis $i$) $\Gamma\vdash A_i$ for each $1\le i\le k$;
\item ($\to$-introduction) from a derivation of $\Gamma,A\vdash B$, a derivation of $\Gamma\vdash A\to B$;
\item ($\to$-elimination) from derivations of $\Gamma\vdash A\to B$ and of $\Gamma\vdash A$, a derivation of $\Gamma\vdash B$.
\end{enumerate}
A derivation is a finite tree labelled by rules; two derivations are equal when their trees and labels are equal.
\end{definition}

\begin{definition}\label{def:stlc}
\emph{Terms} are built from variables $x_1,x_2,\dots$ by abstraction $\lambda x{:}A.\,t$ (with $A$ a formula) and application $t\,u$. A \emph{typing context} is a list $\Delta=(x_1{:}A_1,\dots,x_k{:}A_k)$
of distinct variables with formulas. Typing judgements $\Delta\vdash t:A$ are derived by three rules:
\begin{enumerate}[label=(\roman*)]
\item $\Delta\vdash x_i:A_i$;
\item from $\Delta,x_{k+1}{:}A\vdash t:B$, conclude $\Delta\vdash\lambda x_{k+1}{:}A.\,t:A\to B$;
\item from $\Delta\vdash t:A\to B$ and $\Delta\vdash u:A$, conclude $\Delta\vdash t\,u:B$.
\end{enumerate}
\end{definition}

The bound variable of an abstraction is always the next free index $x_{k+1}$; this convention (a cousin of de~Bruijn indices) means we never need to rename variables. Write $|\Delta|$ for the
context $(A_1,\dots,A_k)$ obtained by forgetting the names.

\begin{lemma}[uniqueness of types]\label{lem:uniquetype}
For every $\Delta$ and $t$ there is at most one $A$ with $\Delta\vdash t:A$, and at most one derivation of it.
\end{lemma}

\begin{proof}
Induction on $t$. The form of $t$ determines the last rule. If $t=x_i$, the rule is (i), which needs $x_i$ in $\Delta$, and the type is $A_i$. If $t=\lambda x_{k+1}{:}A.\,s$, the rule is (ii) and its premise
is a derivation of $\Delta,x_{k+1}{:}A\vdash s:B$; by induction $B$ and that derivation are unique, so the type $A\to B$ and the derivation are. If $t=s\,u$, the rule is (iii), with premises
$\Delta\vdash s:C\to B$ and $\Delta\vdash u:C$ for some $C$; by induction $u$ has at most one type, so $C$ is determined, then $s$ has the unique type $C\to B$, which determines $B$, and both
premise derivations are unique.
\end{proof}

\begin{theorem}[Curry--Howard, implicational fragment]\label{thm:curryhoward}
Let $\Delta=(x_1{:}A_1,\dots,x_k{:}A_k)$ and $A$ be given. Erasing terms is a bijection
\[
\{t:\ \Delta\vdash t:A\}\ \longrightarrow\ \{\text{derivations of }|\Delta|\vdash A\},
\]
from the terms of type $A$ in $\Delta$ to the natural deduction derivations of $A$ from $|\Delta|$. In particular $A$ is derivable from $|\Delta|$ if and only if some term has type $A$ in $\Delta$.
\end{theorem}

\begin{proof}
\emph{The map.} By Lemma~\ref{lem:uniquetype}, a term $t$ with $\Delta\vdash t:A$ has exactly one typing derivation. Erase it: replace each judgement $\Delta'\vdash s:B$ by $|\Delta'|\vdash B$, and each rule
by the rule with the same number (rule (i) for $x_i$ becomes hypothesis $i$). The three typing rules become exactly the three rules of Definition~\ref{def:nd}, so the result is a derivation $\Phi(t)$ of $|\Delta|\vdash A$.

\emph{The inverse.} Given a derivation $D$ of $|\Delta|\vdash A$, define a term $\Psi(D)$ by induction on $D$: if $D$ is hypothesis $i$, $\Psi(D)=x_i$; if $D$ ends with $\to$-introduction from $D'$ deriving
$|\Delta|,B\vdash C$, then $\Psi(D)=\lambda x_{k+1}{:}B.\,\Psi(D')$, where $\Psi(D')$ is computed in the context $\Delta,x_{k+1}{:}B$; if $D$ ends with $\to$-elimination from $D_1$ and $D_2$, then $\Psi(D)=\Psi(D_1)\,\Psi(D_2)$.
By the same induction, $\Delta\vdash\Psi(D):A$, since each case is the matching typing rule.

\emph{They are inverse.} $\Phi(\Psi(D))=D$ by induction on $D$: the typing derivation of $\Psi(D)$ built in the previous step erases, rule by rule, to $D$ (and it is \emph{the} typing derivation of $\Psi(D)$, by uniqueness). Conversely
$\Psi(\Phi(t))=t$ by induction on $t$: if $t=x_i$, $\Phi(t)$ is hypothesis $i$ and $\Psi$ returns $x_i$; if $t=\lambda x_{k+1}{:}B.\,s$, $\Phi(t)$ ends with $\to$-introduction from $\Phi(s)$, and $\Psi$ returns
$\lambda x_{k+1}{:}B.\,\Psi(\Phi(s))=t$; applications are the same.
\end{proof}

So a proof checker for this logic is a type checker for $\lambda$-terms, and the type checker is a dozen lines (Exercise~\ref{ex:typecheck}). Lean's kernel is the same idea for a much richer type
theory, the calculus of inductive constructions, in which types can depend on terms and quantifiers are dependent function types.

\section{What cannot be proved, and why Lean has an axiom}

The logic of Definition~\ref{def:nd} is \emph{intuitionistic}: there is no rule for proof by contradiction. Is anything missing? Peirce's law $((A\to B)\to A)\to A$ is a tautology of classical logic
(check the truth table), but it has no proof in our system. To show that something is \emph{not} derivable we need a semantics in which every derivable formula is true but Peirce's law is not.

\begin{definition}\label{def:heyting3}
Let $H=\{0,\frac12,1\}$ with its usual order, and for $a,b\in H$ put $a\Rightarrow b=1$ if $a\le b$ and $a\Rightarrow b=b$ otherwise. A \emph{valuation} $v$ assigns an element of $H$ to each atom, and extends
to formulas by $v(A\to B)=v(A)\Rightarrow v(B)$. For a context $\Gamma=(A_1,\dots,A_k)$ put $v(\Gamma)=\min(v(A_1),\dots,v(A_k))$, with $v(\Gamma)=1$ for the empty context.
\end{definition}

\begin{lemma}\label{lem:heytingadj}
For all $a,b,g\in H$: $\min(g,a)\le b$ if and only if $g\le(a\Rightarrow b)$.
\end{lemma}

\begin{proof}
If $a\le b$, both sides hold: $\min(g,a)\le a\le b$, and $a\Rightarrow b=1$. If $a>b$, then $a\Rightarrow b=b$, and $\min(g,a)\le b$ holds if and only if $g\le b$ (when $g\le b<a$ the minimum is $g$; when $g>b$ the minimum is $\min(g,a)>b$).
\end{proof}

\begin{theorem}[soundness]\label{thm:heytingsound}
If $\Gamma\vdash A$ is derivable, then $v(\Gamma)\le v(A)$ for every valuation $v$.
\end{theorem}

\begin{proof}
Induction on the derivation. Hypothesis $i$: $v(\Gamma)\le v(A_i)$ as a minimum. Introduction: by induction $\min(v(\Gamma),v(A))=v(\Gamma,A)\le v(B)$, so $v(\Gamma)\le v(A)\Rightarrow v(B)=v(A\to B)$ by Lemma~\ref{lem:heytingadj}.
Elimination: by induction $v(\Gamma)\le v(A)\Rightarrow v(B)$ and $v(\Gamma)\le v(A)$; the lemma with $g=v(\Gamma)$ gives $\min(v(\Gamma),v(A))\le v(B)$, and the left side is $v(\Gamma)$.
\end{proof}

\begin{corollary}\label{cor:peirce}
Peirce's law $((p\to q)\to p)\to p$ is not derivable from the empty context. Consequently there is no closed $\lambda$-term of type $((p\to q)\to p)\to p$.
\end{corollary}

\begin{proof}
Take $v(p)=\frac12$, $v(q)=0$. Then $v(p\to q)=\frac12\Rightarrow0=0$, $v((p\to q)\to p)=0\Rightarrow\frac12=1$, and $v(\text{Peirce})=1\Rightarrow\frac12=\frac12<1=v(\emptyset)$. By Theorem~\ref{thm:heytingsound}
it is not derivable, and by Theorem~\ref{thm:curryhoward} no term has that type.
\end{proof}

This is not a curiosity. In programming terms, Peirce's law is the type of the control operator \texttt{call/cc}, which jumps out of a computation: classical reasoning needs more than pure
functions. Lean's logic is intuitionistic at its core, and classical mathematics is obtained by adding axioms, among them \texttt{Classical.choice}, from which the law of the excluded middle
\texttt{Classical.em} is derived. Mathlib uses them freely; the command \texttt{\#print axioms euclid} lists exactly which axioms a given theorem depends on, which is a kind of honesty
no paper offers.

\begin{think}
The proof of Corollary~\ref{cor:peirce} needed the third value $\frac12$. Why can no two-element semantics show that Peirce's law is not derivable?
\end{think}

\begin{beauty}
Two definitions written for different purposes, one for proofs and one for programs, turn out to be the same definition. Every theorem of this series is, in principle, a program, and
the referee of last resort is a type checker short enough to read in an afternoon. What remains human is the part no kernel checks: choosing what is worth proving, and seeing why it is true.
\end{beauty}

\section{Exercises}

\begin{exercise}[$\star$]\label{ex:leanrun}
Paste the listings of this chapter into \url{https://live.lean-lang.org} and make them compile with the current Mathlib. For each error, say which line of the human proof it corresponds to.
\end{exercise}

\begin{exercise}[$\star$]\label{ex:sterm}
Find a closed term of type $(A\to B\to C)\to(A\to B)\to A\to C$ and draw the natural deduction derivation that Theorem~\ref{thm:curryhoward} associates with it. Then write it in Lean as \texttt{fun f g a =>...}.
\end{exercise}

\begin{exercise}[$\star\star$]\label{ex:typecheck}
Write a type checker for Definition~\ref{def:stlc} in Python: a function that takes a context and a term and returns its type or reports an error. Use Lemma~\ref{lem:uniquetype} to see why it never needs to search.
\end{exercise}

\begin{exercise}[$\star\star$]\label{ex:leanmod9}
Prove in Lean that $x^3+y^3+z^3\not\equiv4\pmod9$ for all \texttt{x y z : ZMod 9}, and deduce the statement for integers by casting with \texttt{Int.cast}. (Expect \texttt{decide} to be slow on $729$ cases; there is a smarter route through \texttt{cubes\_mod\_nine}.)
\end{exercise}

\begin{exercise}[$\star\star$]\label{ex:dne}
Add an atom $\bot$ with $v(\bot)=0$ and put $\neg A=A\to\bot$. Show with Theorem~\ref{thm:heytingsound} that $\neg\neg p\to p$ is not derivable, but find a term of type $p\to\neg\neg p$ and one of type $\neg\neg\neg p\to\neg p$.
\end{exercise}

\begin{exercise}[$\star\star\star$]\label{ex:twovalue}
Answer the question of the \emph{think} box: show that every formula of the implicational fragment that is not derivable takes a value $<1$ under \emph{some} valuation in \emph{some} finite Heyting algebra, but that
the two-element algebra $\{0,1\}$ cannot refute Peirce's law. (The first half is a completeness theorem; you may restrict to formulas with at most two atoms and argue by computer.)
\end{exercise}

\section{Hints}

\begin{hint}{ex:leanrun}
Library names are the usual culprit: search the Mathlib documentation for the statement you need, or use the tactic \texttt{exact?}, which searches the library for you.
\end{hint}

\begin{hint}{ex:sterm}
The answer applies \texttt{f} to \texttt{a} and to \texttt{g a}.
\end{hint}

\begin{hint}{ex:typecheck}
Represent formulas as nested tuples and terms as tagged tuples; recurse on the term.
\end{hint}

\begin{hint}{ex:leanmod9}
For the integers, the cast of $x^3+y^3+z^3$ to \texttt{ZMod 9} is the sum of the cubes of the casts, by \texttt{push\_cast}; and an integer is $\equiv4\pmod9$ exactly when its cast is $4$.
\end{hint}

\begin{hint}{ex:dne}
The valuation $v(p)=\frac12$ works again. For $p\to\neg\neg p$: given $a$ and $f:\neg p$, apply $f$ to $a$.
\end{hint}

\begin{hint}{ex:twovalue}
In $\{0,1\}$ the operation $\Rightarrow$ is classical implication, and Peirce's law is a classical tautology.
\end{hint}

\section{Solutions}

\begin{solution}{ex:leanrun}
\idea A proof assistant is a strict but helpful reader; the error messages point to the gap.
\writeup The typical errors are a renamed lemma (search the documentation or use \texttt{exact?}), a changed implicit argument (add it explicitly with \texttt{(n := n)}), or a tactic that no longer closes a goal
(look at the goal state in the editor and add the missing step). Each is a step a human proof leaves implicit: for instance \texttt{h.le} is ``$p<n$ implies $p\le n$''.
\end{solution}

\begin{solution}{ex:sterm}
\idea Follow the types: to produce $C$ from $f:A\to B\to C$ we need an $A$ and a $B$.
\writeup The term is $\lambda f{:}(A\to B\to C).\,\lambda g{:}(A\to B).\,\lambda a{:}A.\,f\,a\,(g\,a)$, in Lean \texttt{fun f g a => f a (g a)}. Its derivation, in the context $(A\to B\to C,\ A\to B,\ A)$, has three
hypothesis leaves at the top: $\to$-elimination of hypotheses 1 and 3 gives $B\to C$, elimination of hypotheses 2 and 3 gives $B$, one more elimination gives $C$, and three introductions discharge the hypotheses in
reverse order. This is the combinator $S$ of Curry's combinatory logic, and $k$ of the chapter is his $K$.
\end{solution}

\begin{solution}{ex:typecheck}
\idea Uniqueness of types means the checker can compute types bottom-up, deterministically.
\writeup With formulas \texttt{'p'} or \texttt{('->', A, B)} and terms \texttt{('var', i)}, \texttt{('lam', A, t)}, \texttt{('app', t, u)}: for a variable, return the $i$-th type of the context; for \texttt{lam}, check $t$ in the context
extended by $A$ and return \texttt{('->', A, B)}; for \texttt{app}, compute the types of $t$ and $u$, check that the first is \texttt{('->', C, B)} with $C$ equal to the second, and return $B$. Any mismatch is an error. Each case
is one case of the proof of Lemma~\ref{lem:uniquetype}, and the running time is linear in the size of the term (up to comparing formulas). This gives a direct implementation.
\end{solution}

\begin{solution}{ex:leanmod9}
\idea Do the arithmetic where it is finite, then transport.
\writeup Over \texttt{ZMod 9}, \texttt{decide} on all $729$ triples works but is slow; better, by \texttt{cubes\_mod\_nine} each cube is one of $0,1,-1$, and the $27$ combinations (\texttt{rcases} on the three disjunctions, then
\texttt{decide} or \texttt{norm\_num} in each case) give sums in $\{-3,\dots,3\}$, never $4$. For integers, if $x^3+y^3+z^3\equiv4\pmod 9$, cast to \texttt{ZMod 9} (\texttt{push\_cast} moves the cast inside the sum and the powers) to
get a sum of three cubes equal to $4$ there, a contradiction. The equivalence between ``$\equiv4\pmod9$'' and ``cast equals $4$'' is the library lemma \texttt{ZMod.intCast\_eq\_intCast\_iff} or one of its relatives.
\end{solution}

\begin{solution}{ex:dne}
\idea The third value again separates intuitionistic from classical logic.
\writeup With $v(p)=\frac12$: $v(\neg p)=\frac12\Rightarrow0=0$, $v(\neg\neg p)=0\Rightarrow0=1$, and $v(\neg\neg p\to p)=1\Rightarrow\frac12=\frac12<1$; by soundness (the proof of Theorem~\ref{thm:heytingsound} is unchanged by the new atom)
the formula is not derivable. On the other hand $\lambda a{:}p.\,\lambda f{:}\neg p.\,f\,a$ has type $p\to\neg\neg p$, and $\lambda h{:}\neg\neg\neg p.\,\lambda a{:}p.\,h\,(\lambda f{:}\neg p.\,f\,a)$ has type $\neg\neg\neg p\to\neg p$: three negations
reduce to one intuitionistically, though two do not reduce to none.
\end{solution}

\begin{solution}{ex:twovalue}
\idea Two values give classical logic; refuting a classical tautology needs more.
\writeup In $\{0,1\}$, $a\Rightarrow b$ is $0$ exactly when $a=1$ and $b=0$, which is the truth table of implication, so every valuation gives a classical tautology the value $1$; Peirce's law is one, so $\{0,1\}$ cannot refute it, although
by Corollary~\ref{cor:peirce} it is not derivable. The completeness half (the finite model property of intuitionistic logic) goes back to Ja\'skowski in 1936; see S{\o}rensen and Urzyczyn; for formulas with at most two atoms one can check by computer that the non-derivable
ones, found by a proof search that terminates in this fragment, are all refuted in a small finite algebra. We do not prove the general theorem here.
\end{solution}

\section{Sources}

\begin{itemize}
\item T.~C.~Hales et al., A formal proof of the Kepler conjecture, \emph{Forum of Mathematics, Pi} 5 (2017), e2; the introduction recounts the refereeing of the 2005 \emph{Annals} paper and the ``99\% certain'' report.
\item G.~Gonthier, Formal proof---the four-color theorem, \emph{Notices of the AMS} 55 (2008), 1382--1393.
\item P.~Scholze, Liquid tensor experiment, post on the Xena blog, December 2020, and \emph{Experimental Mathematics} 31 (2022); project page \url{https://github.com/leanprover-community/lean-liquid} (completion on 14 July 2022).
\item T.~Tao, blog posts of November and December 2023 on the formalisation of the polynomial Freiman--Ruzsa conjecture; W.~T.~Gowers, B.~Green, F.~Manners, T.~Tao, On a conjecture of Marton, arXiv:2311.05762 (2023).
\item K.~Buzzard, announcement of the EPSRC-funded project to formalise Fermat's Last Theorem (October 2023, project start October 2024), Xena blog.
\item W.~A.~Howard, The formulae-as-types notion of construction, in \emph{To H.\,B.~Curry: Essays on Combinatory Logic, Lambda Calculus and Formalism}, Academic Press, 1980 (notes circulated 1969).
\item M.~H.~S{\o}rensen and P.~Urzyczyn, \emph{Lectures on the Curry--Howard Isomorphism}, Elsevier, 2006: the reference for this section, including Heyting semantics and completeness.
\item The Lean~4 and Mathlib documentation, \url{https://lean-lang.org} and \url{https://leanprover-community.github.io}.
\end{itemize}
